\MethodStyle
\PartDivider{PART IV}{CATRA Part A: Channel Access Control}{Turn local channel observations into station-level contention-window control.}

\begin{frame}{CATRA: Two-Level Fairness Control}
\begin{center}
\resizebox{0.96\textwidth}{!}{%
\begin{tikzpicture}[node distance=0.58cm]
\node[draw=MethodColor,rounded corners,fill=CardBg] (obs) {MAC/PHY observation};
\node[draw=MethodColor,rounded corners,fill=CardBg,right=of obs] (util) {Channel utilization};
\node[draw=MethodColor,rounded corners,fill=CardBg,right=of util] (station) {Station fairness};
\node[draw=MethodColor,rounded corners,fill=CardBg,right=of station] (cw) {Adjust CW};
\draw[-{Latex[length=2mm]},draw=MethodColor] (obs)--(util);
\draw[-{Latex[length=2mm]},draw=MethodColor] (util)--(station);
\draw[-{Latex[length=2mm]},draw=MethodColor] (station)--(cw);
\node[draw=SecondaryColor,rounded corners,fill=CardBg,below=0.9cm of station] (flow) {Flow fairness $\rightarrow$ TCP adaptation};
\draw[-{Latex[length=2mm]},draw=SecondaryColor,dashed] (util)--(flow);
\end{tikzpicture}
}
\end{center}
\vspace{0.55cm}
\KeyPoint{Part A asks: how much channel resource should this station receive, and how much is it actually receiving?}
\end{frame}

\begin{frame}{Ideal Fairness Assumption}
\PaperClaim{In the saturated state, the paper defines ideal per-flow fairness as equal throughput for every flow.}
\begin{itemize}
\item This is a modeling assumption, not a universal definition of network fairness.
\item A station that forwards two flows should receive more channel resource than a station carrying one flow.
\item CATRA translates this flow count into a station-level Fair Bandwidth Ratio.
\end{itemize}
\vspace{0.35cm}
\begin{center}
\begin{tikzpicture}[node distance=1.2cm]
\node[draw=MethodColor,circle,fill=CardBg] (s1) {S1};
\node[below=0.25cm of s1,font=\fontsize{9}{10.5}\selectfont] {2 SEND flows};
\node[draw=MethodColor,circle,fill=CardBg,right=3.0cm of s1] (s2) {S2};
\node[below=0.25cm of s2,font=\fontsize{9}{10.5}\selectfont] {1 SEND flow};
\node[draw=MethodColor,rounded corners,fill=CardBg,right=3.0cm of s2] (share) {$FBR_{S1}:FBR_{S2}=2:1$};
\draw[-{Latex[length=2mm]}] (s1)--(share);
\draw[-{Latex[length=2mm]}] (s2)--(share);
\end{tikzpicture}
\end{center}
\end{frame}

\begin{frame}{SEND, TX, and CS Flows}
\begin{columns}[T,onlytextwidth]
\column{0.32\textwidth}
\textbf{SEND flows}
\begin{itemize}
\item Direct flow generated locally.
\item Forwarding flows transmitted by the examined station.
\end{itemize}
\column{0.32\textwidth}
\textbf{TX flows}
\begin{itemize}
\item Flows from stations in transmission range.
\item Includes locally generated flows.
\end{itemize}
\column{0.32\textwidth}
\textbf{CS flows}
\begin{itemize}
\item Outside transmission range.
\item Inside carrier-sensing range.
\end{itemize}
\end{columns}
\vspace{0.45cm}
\[
n_{total}=n_{TX}+n_{CS},\qquad
n_{CS}=\begin{cases}0,&\text{no CS flow},\\1,&\text{one or more CS flows.}\end{cases}
\]
\end{frame}

\begin{frame}{Station Fair Bandwidth Ratio}
\PaperClaim{The paper defines the station's fair share from the number of flows it sends relative to all flows using its observed channel.}
\[
\boxed{FBR_S=\frac{n_{SEND}}{n_{total}}}
\]
\begin{columns}[T,onlytextwidth]
\column{0.48\textwidth}
\textbf{$n_{SEND}$}
\begin{itemize}
\item Work the station must transmit.
\item Direct plus forwarding flows.
\end{itemize}
\column{0.48\textwidth}
\textbf{$n_{total}$}
\begin{itemize}
\item TX flows plus the CS-flow approximation.
\item Local estimate of competing demand.
\end{itemize}
\end{columns}
\vspace{0.25cm}
\KeyPoint{$FBR_S$ means SHOULD receive; it is not a measured utilization.}
\end{frame}

\begin{frame}{Worked Example: Table 1}
\begin{columns}[T,onlytextwidth]
\column{0.38\textwidth}
\begin{center}
\begin{tikzpicture}[node distance=1.0cm]
\node[draw=MethodColor,circle,fill=CardBg] (s2) {S2};
\node[draw=MethodColor,circle,fill=CardBg,right=of s2] (s1) {S1};
\node[draw=MethodColor,circle,fill=CardBg,right=of s1] (r) {R};
\draw[-{Latex[length=2mm]},draw=MethodColor] (s2)--node[above,font=\fontsize{9}{10.5}\selectfont]{Flow 2}(s1);
\draw[-{Latex[length=2mm]},draw=MethodColor] (s1)--node[above,font=\fontsize{9}{10.5}\selectfont]{Flows 1+2}(r);
\end{tikzpicture}
\end{center}
\column{0.60\textwidth}
\begin{tabular}{lccccc}
\toprule
 & $n_S$ & $n_{TX}$ & $n_{CS}$ & $n_{total}$ & $FBR_S$\\
\midrule
S1 & 2 & 3 & 0 & 3 & $2/3$\\
S2 & 1 & 3 & 0 & 3 & $1/3$\\
R  & 0 & 2 & 1 & 3 & $0$\\
\bottomrule
\end{tabular}
\end{columns}
\vspace{0.55cm}
\begin{itemize}
\item S1 forwards Flow 2 and transmits Flow 1, so it needs twice S2's share.
\item R has no SEND flow; the CS-flow approximation still records competing activity in its view.
\end{itemize}
\end{frame}

\begin{frame}{Real Bandwidth Ratio}
\small
\PaperClaim{CATRA measures the station's active transmission time during an Estimation Period.}
\[
\boxed{RBR_S=\frac{T^{Active}}{EP}}
\]
\begin{columns}[T,onlytextwidth]
\column{0.48\textwidth}
\textbf{$T^{Active}$}
\begin{itemize}
\item Smoothed estimate of airtime used by the station.
\item Includes TCP DATA and TCP ACK transmission activity.
\end{itemize}
\column{0.48\textwidth}
\textbf{$EP$}
\begin{itemize}
\item Fixed observation interval.
\item Original evaluation later selects $EP=2$ s.
\end{itemize}
\end{columns}
\vspace{0.05cm}
\KeyPoint{$FBR_S$ means SHOULD receive; $RBR_S$ means ACTUALLY receiving.}
\end{frame}

\begin{frame}{Algorithm 1: Active-Time Estimation}
\begin{center}
\resizebox{0.96\textwidth}{!}{%
\begin{tikzpicture}[node distance=0.48cm]
\node[draw=MethodColor,rounded corners,fill=CardBg] (observe) {Observe packet $p$};
\node[draw=MethodColor,rounded corners,fill=CardBg,right=of observe] (local) {$destID=localID$?};
\node[draw=MethodColor,rounded corners,fill=CardBg,right=of local] (class) {Classify TCP DATA/ACK};
\node[draw=MethodColor,rounded corners,fill=CardBg,right=of class] (air) {Estimate airtime};
\node[draw=MethodColor,rounded corners,fill=CardBg,right=of air] (sum) {$t\leftarrow t+T_p$};
\draw[-{Latex[length=2mm]}] (observe)--(local);
\draw[-{Latex[length=2mm]}] (local)--(class);
\draw[-{Latex[length=2mm]}] (class)--(air);
\draw[-{Latex[length=2mm]}] (air)--(sum);
\end{tikzpicture}
}
\end{center}
\vspace{0.45cm}
\begin{itemize}
\item Initialize $T^{Active}=0$ and raw accumulator $t=0$.
\item For each locally received packet, inspect MAC and TCP headers.
\item Add the estimated transmission time for TCP DATA or TCP ACK.
\item At every $EP$, smooth the raw observation and reset $t$.
\end{itemize}
\end{frame}

\begin{frame}{Algorithm 1: Packet Airtime Model}
\fontsize{9}{10.8}\selectfont
\[
T_p\approx 0.5\,CW\,ST+T_{RTS}+SIFS+T_{CTS}+SIFS+T_{TCP}+SIFS+T_{ACK}+DIFS
\]
\begin{tabular}{p{3.2cm}p{8.2cm}}
\toprule
\textbf{Term} & \textbf{Interpretation}\\
\midrule
$0.5\,CW\,ST$ & Expected random-backoff time using the current CW.\\
$T_{RTS},T_{CTS}$ & Optional channel-reservation control frames.\\
$SIFS,DIFS$ & Required inter-frame timing.\\
$T_{TCP}$ & Airtime of the TCP DATA or TCP ACK frame.\\
$T_{ACK}$ & Link-layer acknowledgement airtime.\\
\bottomrule
\end{tabular}
\vspace{0.22cm}
\SmallNote{Interpretation: this is an analytical airtime estimate used by CATRA, not a direct PHY-state duration measurement.}
\end{frame}

\begin{frame}{Algorithm 1: Smoothing and Reset}
\[
\boxed{T^{Active}\leftarrow0.8T^{Active}+0.2t},\qquad
\boxed{t\leftarrow0}
\]
\begin{center}
\begin{tikzpicture}[node distance=0.65cm]
\node[draw=MethodColor,rounded corners,fill=CardBg] (old) {80\% previous estimate};
\node[draw=MethodColor,rounded corners,fill=CardBg,right=of old] (new) {20\% current interval};
\node[draw=MethodColor,rounded corners,fill=CardBg,right=of new] (smooth) {Smoothed $T^{Active}$};
\draw[-{Latex[length=2mm]}] (old)--(smooth);
\draw[-{Latex[length=2mm]}] (new)--(smooth);
\end{tikzpicture}
\end{center}
\vspace{0.45cm}
\begin{itemize}
\item Smoothing reduces abrupt control changes caused by one observation period.
\item Resetting $t$ preserves the long-lived estimate while starting a fresh raw measurement.
\item The fixed weights encode a stability/responsiveness trade-off.
\end{itemize}
\end{frame}

\begin{frame}{Adaptive Contention-Window Equation}
\PaperClaim{The station compares actual utilization with its calculated fair share.}
\[
\boxed{CW'=\min\!\left(\frac{RBR_S}{FBR_S}CW,\;CW_{max}\right)}
\]
\begin{columns}[T,onlytextwidth]
\column{0.48\textwidth}
\textbf{Under-allocated: $RBR_S<FBR_S$}
\begin{itemize}
\item $CW'<CW$.
\item Expected backoff decreases.
\item Access opportunity increases.
\end{itemize}
\column{0.48\textwidth}
\textbf{Over-allocated: $RBR_S>FBR_S$}
\begin{itemize}
\item $CW'>CW$ up to $CW_{max}$.
\item Expected backoff increases.
\item Other stations gain opportunities.
\end{itemize}
\end{columns}
\end{frame}

\begin{frame}{CATRA Part A: Complete Feedback Loop}
\begin{center}
\resizebox{0.94\textwidth}{!}{%
\begin{tikzpicture}[node distance=0.55cm]
\node[draw=MethodColor,rounded corners,fill=CardBg] (count) {Count flows};
\node[draw=MethodColor,rounded corners,fill=CardBg,right=of count] (fbr) {$FBR_S$};
\node[draw=MethodColor,rounded corners,fill=CardBg,right=of fbr] (air) {Estimate airtime};
\node[draw=MethodColor,rounded corners,fill=CardBg,right=of air] (rbr) {$RBR_S$};
\node[draw=MethodColor,rounded corners,fill=CardBg,right=of rbr] (ratio) {$RBR_S/FBR_S$};
\node[draw=ConclusionColor,rounded corners,fill=CardBg,right=of ratio] (cw) {Adjust CW};
\draw[-{Latex[length=2mm]}] (count)--(fbr);
\draw[-{Latex[length=2mm]}] (fbr)--(ratio);
\draw[-{Latex[length=2mm]}] (air)--(rbr);
\draw[-{Latex[length=2mm]}] (rbr)--(ratio);
\draw[-{Latex[length=2mm]}] (ratio)--(cw);
\end{tikzpicture}
}
\end{center}
\vspace{0.55cm}
\begin{itemize}
\item Flow counts define the target share.
\item Airtime estimation measures the realized share.
\item CW is the MAC-layer actuator.
\item The result targets per-station fairness, not yet per-flow fairness.
\end{itemize}
\end{frame}
